\subsection{Method}
Applied \texttt{HallucinationDetector.detect\_all} (spectral/structural/textural) to EDSR-small and bicubic SR outputs vs HR reference on 1,742 clean test tiles. Semantic detection skipped (no task model). Resume-safe per-tile JSONL.

\subsection{Results}

\begin{table}[h]
\centering
\caption{Phase~6 hallucination detection rates on 1,742 clean test tiles.}
\label{tab:hall}
\begin{tabular}{lrrr}
\toprule
Category & EDSR-small & Bicubic $\times$4 & McNemar $p$ \\
\midrule
Spectral & 2.53\% (44) & 3.62\% (63) & $3.6 \times 10^{-5}$ \\
Structural & 0.11\% (2) & 0.00\% (0) & 0.48 (n.s.) \\
Textural & 4.54\% (79) & 9.82\% (171) & $< 0.001$ \\
\textbf{Any} & \textbf{6.89\% (120)} & \textbf{13.03\% (226)} & $< 0.001$ \\
\bottomrule
\end{tabular}
\end{table}

\begin{itemize}
    \item EDSR produces \textbf{half} the hallucination rate of bicubic ($-$6.14 pp)
    \item Textural hallucinations dominate and show the largest gap ($-$5.28 pp)
    \item Spectral hallucinations rare for both (respects band correlations)
    \item Structural negligible for both
\end{itemize}

\subsection{Key Finding}
EDSR-small improves \textbf{both} fidelity (+3.03 dB PSNR) \textbf{and} reduces hallucinations (6.89\% vs 13.03\%). This contradicts the common assumption that higher PSNR implies more hallucination -- the learned model generates textures that are more plausible AND more faithful to the HR reference.
